% ======================================================================
% Tobacco Smoke Remediation Practices Among Residential Cleaning and
% Restoration Professionals in San Diego County
% Amanda Ta - PSY 499: Special Study (San Diego State University)
%
% LaTeX edition of the author's course report. The prose is unchanged
% from the original Word manuscript; only the typesetting is new.
%
% To compile: pdflatex Ta_PSY499_Tobacco_Smoke_Remediation.tex
% Requires the apa7 document class (TeX Live: tlmgr install apa7).
% ======================================================================

\documentclass[stu]{apa7}

\usepackage[utf8]{inputenc}
\usepackage{hyperref}    % clickable URLs

\title{Tobacco Smoke Remediation Practices Among Residential Cleaning and Restoration Professionals in San Diego County}
\shorttitle{Tobacco Smoke Remediation Practices}
\author{Amanda Ta}
\affiliation{School of Public Health, San Diego State University}
\course{PSY 499: Special Study}
\professor{Dr. Georg E. Matt, PhD}
\duedate{May 21, 2026}

\begin{document}
\maketitle

\begin{center}
\textbf{Tobacco Smoke Remediation Practices Among Residential Cleaning and Restoration Professionals in San Diego County}
\end{center}

Commercial tobacco smoke contamination remains a persistent indoor environmental health concern, particularly in residential settings where tobacco smoke residue can accumulate on indoor surfaces, furnishings, dust, and ventilation systems long after active smoking has stopped. This residual contamination, commonly referred to as thirdhand smoke (THS) (Matt et al., 2011), contains nicotine, tobacco-specific nitrosamines, polycyclic aromatic hydrocarbons, volatile organic compounds, heavy metals, and other toxic constituents (Matt et al., 2011; U.S. Environmental Protection Agency, 2024) that may persist indoors for extended periods of time and contribute to ongoing exposure among nonsmokers (Matt et al., 2011). Research conducted by the Center for Tobacco and the Environment (CTE) at San Diego State University and related thirdhand smoke studies has demonstrated that THS pollutants can remain detectable in homes previously occupied by smokers even after residents move out and visible signs of smoking have been removed. These pollutants may re-emit into indoor air, accumulate in household dust, and react with other indoor pollutants to form secondary toxic compounds, raising concerns regarding chronic exposure in residential environments (Sleiman et al., 2010).

While growing evidence has documented the persistence and potential health implications of THS contamination, substantially less is known about how remediation and cleaning professionals address tobacco smoke residue in real-world residential settings. Existing remediation practices for tobacco smoke contamination vary widely and may involve combinations of surface cleaning, carpet and upholstery cleaning, HVAC cleaning, air scrubbing, ozone treatment, hydroxyl generation, thermal fogging, repainting, sealing, and removal of contaminated materials (Matt et al., 2011; Winickoff et al., 2009). However, few studies have systematically examined which remediation approaches are currently used in practice, how professionals assess contamination severity, what methods are perceived as most effective, and how remediation decisions are made across different levels of tobacco smoke contamination.

Residential tobacco smoke contamination may present challenges in multiunit housing and rental properties, where contamination from previous occupants can persist despite turnover cleaning or cosmetic renovation (Matt et al., 2011). Concerns regarding odor, staining, and potential health effects may also disproportionately affect lower-income households and renters with limited control over prior indoor smoking conditions (Matt et al., 2011). Prior CTE research has emphasized the need for improved understanding of remediation, and exposure reduction strategies for tobacco smoke contamination in homes and other indoor environments. Understanding current remediation practices may help identify gaps in remediation knowledge, inconsistencies in remediation standards, and opportunities to improve evidence-based guidance for reducing tobacco-related indoor environmental exposure (Matt et al., 2011).

To address these gaps, we conducted a descriptive mixed-methods study of remediation and cleaning professionals operating in San Diego County, California that combined quantitative characterization of remediation-related businesses and professional certifications with qualitative interviews examining how professionals assess and remediate residential tobacco smoke contamination, including remediation workflows, remediation technologies and products used, perceived effectiveness, and estimated service costs across varying contamination scenarios.

\section{Methods}

\subsection{Study Design}

We conducted a descriptive mixed-methods study to examine how remediation and cleaning professionals in San Diego County assess and remediate residential tobacco smoke contamination using quantitative characterization of remediation-related businesses and professional certifications and semi-structured qualitative interviews regarding remediation practices, perceived effectiveness, and estimated service costs.

The study focused on remediation methods, technologies and products used, perceived remediation effectiveness, and estimated service costs under varying contamination scenarios. This work aligns with ongoing research conducted by the Center for Tobacco and the Environment (CTE) at San Diego State University examining environmental tobacco toxicants and remediation strategies in indoor environments (Matt et al., 2011).

The study combined systematic database development with semi-structured interviews. First, we developed and verified a database of remediation-related companies and certified professionals operating within San Diego County. Second, we recruited eligible participants from the database to complete structured telephone or Zoom interviews regarding their professional experience with tobacco smoke remediation.

The qualitative interview component was designed to capture detailed information regarding remediation workflows, assessment procedures, cleaning and restoration technologies, use of odor-control methods, and decision-making processes used in homes affected by tobacco smoke residue. Participants were also asked to estimate remediation costs and timelines for standardized hypothetical residential contamination scenarios. Interviews followed a semi-structured format to ensure consistency across participants while allowing flexibility to explore remediation approaches and professional experiences in greater depth.

The study protocol incorporated standardized recruitment, screening, consent, recording, and interview procedures developed by the research team. All interviews were audio-recorded with participant permission and transcribed for qualitative analysis.

\subsection{Study Population}

The study population consisted of remediation and cleaning professionals operating within San Diego County, California with experience addressing residential tobacco smoke contamination, including cigarette smoke residue, stale cigarette odor, and thirdhand smoke-related contamination. Eligible businesses included restoration and remediation companies, general cleaning services, carpet cleaning companies, and individually certified remediation professionals. We prioritized companies and professionals holding certifications relevant to smoke and odor remediation, including Fire and Smoke Restoration Technician (FSRT), Odor Control Technician (OCT), Carpet Cleaning Technician (CCT), Carpet Repair and Reinstallation Technician (RRT), Upholstery and Fabric Cleaning Technician (UFT), and House Cleaning Technician (HCT) certifications through the Institute of Inspection Cleaning and Restoration Certification (IICRC).

\subsection{Database Development}

We developed a structured database of remediation and cleaning professionals potentially involved in residential tobacco smoke remediation within San Diego County. Database development was conducted in multiple stages using existing laboratory records, professional certification databases, and publicly available business information.

We first reviewed a pre-existing 2020 laboratory dataset containing 24 companies previously identified as potentially relevant to smoke remediation services. Each company was re-evaluated to confirm current business activity, verify contact information, and assess continued relevance to tobacco smoke remediation.

We then conducted an additional search using the Institute of Inspection Cleaning and Restoration Certification (IICRC) Global Locator, a public directory of certified remediation and cleaning professionals. Searches were conducted within a 250-mile radius of San Diego, California to identify companies and individual professionals holding certifications associated with smoke restoration, odor control, carpet cleaning, and related remediation services.

Because the IICRC search radius extended beyond the intended study region, we manually reviewed company websites, public business listings, certification records, and publicly available contact information to verify whether businesses actively operated within San Diego County or surrounding service areas. Companies were excluded if they were inactive, lacked evidence of relevant remediation services, or did not operate within the study area.

Companies were classified into four categories: restoration/remediation companies, general cleaning companies, carpet cleaning companies, and individually certified remediation professionals. For each company or professional, we collected information including company name, business category, website, phone number, email address, city, ZIP code, certification status, certifications held, data source, and notes regarding remediation services and recruitment eligibility. We also documented whether businesses specifically advertised cigarette smoke remediation, odor removal, ozone treatment, thermal fogging, surface cleaning, HVAC cleaning, or related remediation approaches relevant to tobacco smoke contamination.

The final database included remediation and cleaning professionals operating across the San Diego region and served as the sampling frame for participant recruitment and interview procedures.

\subsection{Participant Sampling and Recruitment}

We used purposive sampling to recruit remediation and cleaning professionals with experience addressing residential tobacco smoke contamination. The company database developed for this study served as the primary sampling frame for recruitment activities. Because the study aimed to characterize remediation approaches used in real-world residential settings, we prioritized businesses and professionals with demonstrated experience in smoke restoration, odor control, environmental cleaning, and related remediation services.

To improve representation across remediation-related industries, we recruited participants from multiple professional categories, including restoration/remediation companies, general cleaning services, carpet cleaning companies, and individually certified remediation professionals. We further prioritized companies holding certifications relevant to smoke and odor remediation, particularly Fire and Smoke Restoration Technician (FSRT) and Odor Control Technician (OCT) certifications through the Institute of Inspection Cleaning and Restoration Certification (IICRC).

Companies were assigned randomized identification numbers and organized into prioritized recruitment tiers to reduce systematic recruitment bias while maintaining emphasis on smoke-remediation expertise. Randomized recruitment order was generated using Random.org prior to participant contact. We aimed to recruit participants across multiple service categories, including remediation/restoration, general cleaning, and carpet cleaning professionals.

Recruitment was conducted by trained research staff using telephone outreach procedures. Companies were contacted using publicly available business phone numbers and contact information collected during database development. If a business receptionist or staff member was unable to answer study questions, researchers requested referral to a company representative with direct knowledge of tobacco smoke remediation practices, such as an owner, manager, estimator, or technician.

Participants were informed that the study was being conducted by researchers at San Diego State University's Center for Tobacco and the Environment to better understand remediation approaches used for tobacco smoke residue in homes. Eligible participants were invited to complete a 45--60 minute telephone or Zoom interview regarding remediation methods, technologies, perceived effectiveness, and service costs associated with tobacco smoke contamination. Participants who completed the interview received a \$75 Amazon gift card incentive.

We attempted to contact each company up to three times before classifying the business as nonresponsive or unreachable. In addition to purposive recruitment, we incorporated snowball sampling by asking participants whether they recommended additional professionals or subcontractors with relevant experience in tobacco smoke remediation.

\subsection{Interview Procedures and Measures}

Eligible participants completed semi-structured interviews conducted by trained research staff affiliated with the Center for Tobacco and the Environment (CTE) at San Diego State University. Interviews were conducted by telephone or Zoom according to participant preference and were designed to collect detailed information regarding residential tobacco smoke remediation practices, remediation technologies, perceived effectiveness, and service costs. Prior to participation, eligible participants received information regarding study procedures, confidentiality protections, audio recording, and participant compensation. Participants were offered the option to review the broad consent form electronically before the interview or have the consent form read aloud at the beginning of the interview session.

At the beginning of each interview, research staff reviewed the study purpose, voluntary participation, confidentiality protections, and recording procedures using a standardized introduction script. Participants were informed that the study aimed to better understand how remediation and cleaning professionals assess and remediate residential tobacco smoke contamination, including remediation approaches, technologies, and costs. All interviews were audio-recorded with participant permission for note-taking and transcription purposes. Zoom's automatic transcription feature was used during Zoom interviews to support transcription and qualitative data management procedures. Interviews typically lasted approximately 45--60 minutes.

We developed a semi-structured interview guide informed by prior CTE research on thirdhand smoke contamination and remediation, internal project development meetings, remediation industry practices identified during database construction, and topics identified during recruitment and screening activities. Interview domains included participant background and professional experience, remediation assessment procedures, remediation workflows, remediation technologies and products used, perceived remediation effectiveness, client experiences, and remediation cost estimation.

Participants were asked about the types of services they provide for homes affected by tobacco smoke contamination, including deep cleaning, HVAC cleaning, deodorization, air scrubbing, ozone treatment, thermal fogging, repainting, and removal of contaminated materials. Additional questions examined how participants assessed contamination severity, selected remediation approaches, and modified remediation procedures according to odor intensity, visible staining, or extent of contamination.

The interview guide also included questions regarding specific remediation technologies and products, including air scrubbers, ozone generators, hydroxyl generators, trisodium phosphate (TSP), thermal fogging devices, specialty primers, sealants, and cleaning products commonly used during remediation. Participants described how remediation technologies were used, safety precautions associated with treatment procedures, occupancy restrictions, and factors influencing perceived remediation effectiveness.

Finally, participants completed a cost estimation section involving three standardized hypothetical residential contamination scenarios. Each scenario described a 600-square-foot, one-bedroom apartment with differing smoking histories and contamination severity levels. Participants were asked to estimate recommended remediation approaches, anticipated timelines, and approximate service costs for each scenario based on their professional experience. Interview guides, screening materials, and consent documents were developed for study procedures.

\subsection{Ethical Considerations}

This study was conducted through the Center for Tobacco and the Environment (CTE) at San Diego State University and followed procedures developed for human subjects research involving qualitative interviews and recruitment of adult professionals. Eligible participants were informed that participation was voluntary and that they could decline to answer any question or withdraw from the interview at any time without penalty.

Prior to participation, participants were provided information regarding study procedures, confidentiality protections, audio recording, and participant compensation. Participants were offered the option to review the broad consent form electronically before the interview or have the consent form read aloud at the beginning of the interview session. Broad consent procedures informed participants that anonymized data collected during the study could be retained for future tobacco-related research conducted by the Center for Tobacco and the Environment.

All interviews were audio-recorded only with participant permission. Participants were informed that no identifying information about individuals or businesses would be included in reports, publications, or presentations resulting from the study. Recruitment tracking information, participant contact information, and interview materials were stored separately from interview responses and protected using password-restricted research files accessible only to authorized study personnel.

\section{References}

\noindent\hangindent=0.5in Matt, G. E., Quintana, P. J. E., Destaillats, H., Gundel, L. A., Sleiman, M., Singer, B. C., Jacob, P., Benowitz, N., Winickoff, J. P., Rehan, V., Talbot, P., Schick, S., Samet, J., Wang, Y., Hang, B., Martins-Green, M., Pankow, J. F., \& Hovell, M. F. (2011). Thirdhand tobacco smoke: Emerging evidence and arguments for a multidisciplinary research agenda. \textit{Environmental Health Perspectives}, \textit{119}(9), 1218--1226. https://doi.org/10.1289/ehp.1103500

\noindent\hangindent=0.5in Matt, G. E., Quintana, P. J. E., Hovell, M. F., Bernert, J. T., Song, S., Novianti, N., Juarez, T., Floro, J., Gehrman, C., Garcia, M., \& Larson, S. (2008). Residual tobacco smoke pollution in used cars for sale: Air, dust, and surfaces. \textit{Nicotine \& Tobacco Research}, \textit{10}(9), 1467--1475. https://doi.org/10.1080/14622200802279898

\noindent\hangindent=0.5in Matt, G. E., Quintana, P. J. E., Zakarian, J. M., Fortmann, A. L., Chatfield, D. A., Hoh, E., Uribe, A. M., \& Hovell, M. F. (2011). When smokers move out and non-smokers move in: Residential thirdhand smoke pollution and exposure. \textit{Tobacco Control}, \textit{20}(1), e1. https://doi.org/10.1136/tc.2010.037382

\noindent\hangindent=0.5in Sleiman, M., Gundel, L. A., Pankow, J. F., Jacob, P., Singer, B. C., \& Destaillats, H. (2010). Formation of carcinogens indoors by surface-mediated reactions of nicotine with nitrous acid, leading to potential thirdhand smoke hazards. \textit{Proceedings of the National Academy of Sciences}, \textit{107}(15), 6576--6581. https://doi.org/10.1073/pnas.0912820107

\noindent\hangindent=0.5in U.S. Environmental Protection Agency. (2024). \textit{Volatile organic compounds' impact on indoor air quality.} https://www.epa.gov/indoor-air-quality-iaq/volatile-organic-compounds-impact-indoor-air-quality

\noindent\hangindent=0.5in Winickoff, J. P., Friebely, J., Tanski, S. E., Sherrod, C., Matt, G. E., Hovell, M. F., \& McMillen, R. C. (2009). Beliefs about the health effects of ``thirdhand'' smoke and home smoking bans. \textit{Pediatrics}, \textit{123}(1), e74--e79. https://doi.org/10.1542/peds.2008-2184

\end{document}
